\section{Gabarito e resoluções comentadas — Direito Administrativo}

\begin{solucao}{04.01}\textbf{CERTO.} O sentido subjetivo/orgânico aponta os sujeitos da Administração; o objetivo/material, a atividade administrativa.\end{solucao}
\begin{solucao}{04.02}\textbf{ERRADO.} Distribuir competências dentro da mesma pessoa jurídica é desconcentração. Descentralização envolve outra pessoa jurídica ou particular.\end{solucao}
\begin{solucao}{04.03}\textbf{CERTO.} Órgão é centro de competências sem personalidade jurídica própria.\end{solucao}
\begin{solucao}{04.04}\textbf{ERRADO.} Nem todo ato é imperativo ou autoexecutório; esses atributos dependem da natureza e do regime do ato.\end{solucao}
\begin{solucao}{04.05}\textbf{CERTO.} Anulação trata ilegalidade; revogação, mérito de ato válido, respeitados direitos e limites.\end{solucao}
\begin{solucao}{04.06}\textbf{ERRADO.} O Judiciário controla juridicidade, não substitui a Administração no mérito legítimo para revogar ato válido.\end{solucao}
\begin{solucao}{04.07}\textbf{ERRADO.} Efetividade decorre da investidura em cargo efetivo; estabilidade é garantia adquirida sob requisitos constitucionais próprios.\end{solucao}
\begin{solucao}{04.08}\textbf{CERTO.} Vinculação e supervisão finalística não criam subordinação hierárquica entre pessoas jurídicas distintas.\end{solucao}
\begin{solucao}{04.09}\textbf{CERTO.} Inexigibilidade pressupõe inviabilidade de competição; dispensa decorre de hipótese legal de contratação direta.\end{solucao}
\begin{solucao}{04.10}\textbf{ERRADO.} A reforma de 2021 reforçou o dolo nos tipos de improbidade, afastando a tese de culpa suficiente em qualquer hipótese.\end{solucao}

\begin{solucao}{04.11}\textbf{B.} Criar secretaria interna é desconcentrar competências sem criar nova pessoa jurídica.\end{solucao}
\begin{solucao}{04.12}\textbf{C.} Sociedade de economia mista é entidade com personalidade própria da Administração indireta. Autarquia também é entidade, não órgão.\end{solucao}
\begin{solucao}{04.13}\textbf{CERTO.} Discricionariedade é liberdade juridicamente delimitada, não imunidade ao controle.\end{solucao}
\begin{solucao}{04.14}\textbf{B.} Vício de competência não exclusiva pode ser convalidável, desde que presentes os requisitos legais e ausentes prejuízos impeditivos.\end{solucao}
\begin{solucao}{04.15}\textbf{B.} Ato válido que se torna inconveniente pode ser revogado pela Administração, observados limites jurídicos.\end{solucao}
\begin{solucao}{04.16}\textbf{CERTO.} A presunção é juris tantum, podendo ser afastada por prova.\end{solucao}
\begin{solucao}{04.17}\textbf{B.} O agente tem competência, mas a usa para finalidade imprópria: desvio de finalidade.\end{solucao}
\begin{solucao}{04.18}\textbf{B.} O poder hierárquico organiza e controla internamente relações de subordinação.\end{solucao}
\begin{solucao}{04.19}\textbf{CERTO.} Essa é a distinção clássica entre as modalidades de abuso de poder.\end{solucao}
\begin{solucao}{04.20}\textbf{B.} O art. 37, §6º, adota responsabilidade objetiva da pessoa jurídica, preservado o regresso contra o agente nos casos constitucionais.\end{solucao}
\begin{solucao}{04.21}\textbf{ERRADO.} Responsabilidade objetiva dispensa culpa do Estado, não dano e nexo causal.\end{solucao}
\begin{solucao}{04.22}\textbf{C.} Estabilidade não nasce da posse e não se confunde com efetividade ou vitaliciedade.\end{solucao}
\begin{solucao}{04.23}\textbf{B.} A avocação é excepcional e temporária, devendo ser motivada dentro dos limites legais.\end{solucao}
\begin{solucao}{04.24}\textbf{CERTO.} A Lei nº 9.784/1999 contém hipóteses de competência indelegável, entre elas matérias normativas e decisão de recursos.\end{solucao}
\begin{solucao}{04.25}\textbf{B.} Inviabilidade de competição é o núcleo da inexigibilidade.\end{solucao}

\begin{solucao}{04.26}\textbf{C.} I e II estão corretas. III é falsa porque o Judiciário não revoga ato válido por juízo de conveniência.\end{solucao}
\begin{solucao}{04.27}\textbf{CERTO.} O poder de autotutela convive com decadência e segurança jurídica, ressalvada má-fé nos termos legais.\end{solucao}
\begin{solucao}{04.28}\textbf{B.} A decisão sancionatória precisa expor fundamentos de fato e de direito suficientes.\end{solucao}
\begin{solucao}{04.29}\textbf{B.} Autotutela permite à Administração invalidar seus atos ilegais e rever atos válidos dentro do regime aplicável.\end{solucao}
\begin{solucao}{04.30}\textbf{CERTO.} A invalidação de situação individual pode exigir contraditório e ampla defesa, conforme o contexto e o regime jurídico.\end{solucao}
\begin{solucao}{04.31}\textbf{B.} Continuidade é princípio, mas o ordenamento admite interrupções justificadas e reguladas.\end{solucao}
\begin{solucao}{04.32}\textbf{B.} Omissões exigem análise do dever de agir e do nexo, com atenção ao tratamento jurisprudencial específico.\end{solucao}
\begin{solucao}{04.33}\textbf{CERTO.} Culpa exclusiva da vítima pode romper o nexo; culpa concorrente pode influenciar o resultado conforme o caso.\end{solucao}
\begin{solucao}{04.34}\textbf{B.} Nem toda ilegalidade é improbidade; é necessário verificar o tipo e o elemento subjetivo exigido na lei vigente.\end{solucao}
\begin{solucao}{04.35}\textbf{A.} Preferência administrativa não cria inviabilidade de competição. O processo deve demonstrar objetivamente o enquadramento.\end{solucao}
\begin{solucao}{04.36}\textbf{CERTO.} Contratação direta continua sujeita a instrução e motivação, inclusive quanto a preço e hipótese legal.\end{solucao}
\begin{solucao}{04.37}\textbf{B.} Tribunais de Contas auxiliam o Poder Legislativo no controle externo, conforme desenho constitucional.\end{solucao}
\begin{solucao}{04.38}\textbf{B.} Autarquia é pessoa jurídica própria; há vinculação e supervisão, não relação hierárquica comum com o ministério.\end{solucao}
\begin{solucao}{04.39}\textbf{CERTO.} A teoria dos motivos determinantes torna relevante a correspondência entre os motivos declarados e a realidade, inclusive em atos discricionários.\end{solucao}
\begin{solucao}{04.40}\textbf{B.} As responsabilidades podem coexistir, com independência relativa e efeitos legais entre instâncias.\end{solucao}

\begin{solucao}{04.41}\textbf{CERTO.} O poder de invalidar não é ilimitado no tempo e deve respeitar decadência, boa-fé e garantias processuais aplicáveis.\end{solucao}
\begin{solucao}{04.42}\textbf{B.} Eficiência não autoriza o Judiciário a substituir opção administrativa legítima por preferência própria sem vício jurídico.\end{solucao}
\begin{solucao}{04.43}\textbf{B.} O achado deve demonstrar a ausência do pressuposto da inexigibilidade e indicar análise da instrução e eventual responsabilização, sem generalizar.\end{solucao}
\begin{solucao}{04.44}\textbf{CERTO.} Ilegalidade e improbidade são categorias distintas; a segunda exige requisitos específicos.\end{solucao}
\begin{solucao}{04.45}\textbf{B.} Poder de polícia é prerrogativa juridicamente fundada; falta de competência é vício central.\end{solucao}
\begin{solucao}{04.46}\textbf{CERTO.} Motivos declarados podem vincular a validade do ato, permitindo controle de sua existência e veracidade.\end{solucao}
\begin{solucao}{04.47}\textbf{B.} A Constituição inclui pessoas jurídicas privadas prestadoras de serviço público no regime do art. 37, §6º.\end{solucao}
\begin{solucao}{04.48}\textbf{B.} Decisão que ignora argumentos relevantes pode falhar em motivação e devido processo.\end{solucao}
\begin{solucao}{04.49}\textbf{ERRADO.} Eficiência opera dentro da legalidade; não autoriza afastar exigência normativa.\end{solucao}
\begin{solucao}{04.50}\textbf{B.} A decisão de prorrogar deve ser motivada e baseada em desempenho, necessidade, vantajosidade e alternativas, não em inércia.\end{solucao}
